%! TEX root = main.tex

\clearpage

{\noindent \bfseries Abstract \par}
\vspace{1cm}
Since the emergence of generative AI services, academia and industry worldwide have begun to demand massive amounts of data for model training. As the demand for high-quality data in various fields has surged to unprecedented levels, data has attained a status akin to a strategic asset. Consequently, platform companies that previously monopolized data in their respective fields have exerted immense influence. Meanwhile, the individuals who contributed to the creation of that data have received no compensation and continue to rely on the untrustworthy infrastructure of corporations. Web 3.0 has begun to receive spotlight as a vision to correct ownership structures in this era of data assets, and many engineers have focused on technological development for this purpose. In particular, numerous studies have aimed to implement systems that allow for the free and safe trading of data while protecting the ownership of the original data author. As part of such an effort, this study proposes a system that uses zero-knowledge proof technology to search for 'evidence' rather than the original raw data. Specifically, by proposing a system that can search for numerical medical data based on ranges, individual patients can upload evidence to a platform without fear of their medical data being leaked. Data consumers can then search for this evidence to verify the existence of data that fits their required range. By storing only zk-SNARK-based zero-knowledge proof data on a centralized server, this system eliminates the risk of personal information leakage while simultaneously providing an efficient search mechanism for the data ranges desired by consumers. \\ \\

{\noindent \bfseries Keywords \par}
\noindent data privacy, zero knowledge proof, data exchange platform, zk-SNARK, medical data, range proof
